\scanpage{68}
% archive sheet number: 67
% Two long diagonal strokes cancel this draft. Unfilled sign slots are left as spacing.
\deleted{sous les conditions envisagées a) b) $\quad)$ $\quad)$, $B$ plat sur $A$ et à extension résiduelle finie,}
\[
\deleted{\text{\uuline{R}}(A)\quad\text{et}\quad\text{\uuline{R}}(B\otimes k)\ \text{géométrique}\quad\text{impliquent}\quad\text{\uuline{R}}(B).}
\]
\deleted{La condition supplémentaire qu'on vient de signaler pour \uuline{R} est satisfaite dans tous les exemples envisagés dans 7.5.3., à l'exception de (vii), où il est en défaut. En effet, on peut (même \deleted{sur un corps} pour une courbe $Y$ sur un corps algébriquement clos $k$ de caractéristique arbitraire) donner des exemples d'un morphisme étale $f\colon X\quad Y$, et d'un point $x$ de $X$ tel que $Y$ soit intègre en $y=f(x)$, mais $X$ ne le soit pas en $x$.\handwritten{/} Ce serait peut-être le moment de donner l'exemple (également signalé dans 6.15.7. et qui resert $X$ fois): on prend une courbe à point double ordinaire, et on recolle deux exemplaires du normalisé \unsure{d|de} de façon à avoir une courbe étale $X$ au dessus de $Y$. Un exemple analogue avec un anneau local \inserted{intègre} \uline{complet}, donc unibranche, mais en l'occurrence non géométriquement unibranche, est donné par $\inserted{A=}\text{\uuline{R}}\ X,Y\quad/(X^{\handwritten{2}}+Y^{\handwritten{2}})$. \deleted{Si $X$ est son spectre} Alors $A\otimes_{\text{\uuline{R}}}\underline{C}$ est étale sur $A$, mais non intègre. Ces exemples, qu'on semble rencontrer pour la première fois à propos de 6.15.7. pourraient peut-être remonter à cet endroit.\handwritten{/} Il est cependant plausible que \deleted{l'énoncé qui} la reponse à la question suivante soit affirmative :}

\deleted{Soit $A\quad B$ un homomorphisme local d'anneaux locaux noethériens, $B$ plat sur $A$, $B\otimes_A k\handwritten{/}k$ géométriquement intègre, $A$ \deleted{excellent (par exemple complet)} géométriquement unibranche (NB on peut se ramener au cas: $A$ normal). Si $A$ est de plus excellent (disons complet), alors $B$ est-il intègre ?}

\deleted{La reponse est affirmative quand le corps résiduel $k$ de $A$ est de car 0, comme on verra au N° 7.9. en utilisant la\deleted{\unsure{x}} résolution des singularités de Hironaka, ou lorsque $B$ est essentiellement de type fini sur $A$ (réf $\qquad$). Mais \inserted{dans le cas général,} même lorsque l'extension résiduelle est triviale, $A$\deleted{\unsure{s|x}} et $B$ complet, \deleted{\unsure{la technique dévis}} nous ignorons si la réponse est affirmative. On a cependant un résultat partiel ci-dessous.}
